\section{正合拉格朗日子流形}
\label{sec:exact-lagrangian-submanifolds}

余切丛上的余法丛只是拉格朗日子流形非常特殊的例子；拉格朗日子流形可以定义在一般辛流形上。

在开始一般性研究拉格朗日子流形之前，先看几个例子。

\begin{example}\label{ex:exact-lagrangian-cotangent}
考虑带典范辛形式的余切丛 $M=T^*N$。
\begin{enumerate}[label=(\arabic*)]
\item 对余法丛 $\nu^*S$，有
\begin{equation}
 i^*\theta\equiv0.
 \label{eq:conormal-liouville-zero}
\end{equation}
\item 闭 $1$-形式的图像：设 $\beta$ 是 $N$ 上的闭 $1$-形式，并把 $\beta$ 看作相应的截面映射
\[
 \widetilde\beta:N\longrightarrow T^*N,
\]
其像记作 $\operatorname{Graph}\beta$。在这种情形下，由（2.1.2）可得
\begin{equation}
 d(\widetilde\beta^*\theta)=d\beta=0.
 \label{eq:closed-one-form-graph}
\end{equation}
特别地，光滑函数 $f:N\to\mathbb R$ 的微分 $df$ 的图像是拉格朗日的。
\end{enumerate}
\end{example}

正合 $1$-形式 $df$ 的图像，是所谓正合拉格朗日子流形的特例。

\begin{definition}\label{def:exact-symplectic-manifold}
若辛流形 $(M,\omega)$ 上存在 $1$-形式 $\alpha$ 使 $\omega=d\alpha$，则称它是\emph{正合的}。
\end{definition}

在正合辛流形 $(M,d\alpha)$ 上，$i:L\to M$ 是拉格朗日的，当且仅当
\stfhIndex{I122}{zheng^he^xin^liu^xing}{正合辛流形（exact symplectic manifold）}
$i^*\alpha$ 闭。

\begin{definition}\label{def:exact-lagrangian-submanifold}
正合辛流形 $(M,d\alpha)$ 中的拉格朗日子流形 $i:L\to M$ 称为\emph{正合的}，如果 $i^*\alpha$ 正合。
\stfhIndex{I121}{zheng^he^la^ge^lang^ri^zi^liu^xing}{正合拉格朗日子流形（exact Lagrangian submanifold）}
\end{definition}

正合性定义立即蕴含：存在函数 $f:L\to\mathbb R$ 使 $i^*\alpha=df$。$f$ 的选取在相差加法常数的意义下唯一。稍后会看到，在作用量泛函的变分研究中，把满足 $df=i^*\alpha$ 的配对 $(L,f)$ 看作正合拉格朗日子流形十分重要。

正合拉格朗日子流形的另一个基本性质，是下面这个\emph{周期}消失结论；它是斯托克斯定理的直接推论。

\begin{proposition}\label{prop:exact-lagrangian-period-vanishing}
设 $L\subset(M,\omega)$ 是正合拉格朗日子流形，$\omega=d\alpha$。设
$w:(\Sigma,\partial\Sigma)\to(M,L)$ 是从带边界紧二维曲面 $\Sigma$ 出发的任意光滑映射。则
\[
 \int w^*\omega=0.
\]
\end{proposition}

\begin{proof}
利用 $\omega=d\alpha$ 并应用斯托克斯公式，得到
\[
 \int_\Sigma w^*\omega=\int_\Sigma w^*d\alpha
 =\int_{\partial\Sigma}w^*\alpha.
\]
由 $L$ 的正合性，对 $L$ 上某个函数 $f$ 有 $i_L^*\alpha=df$。因此
$w^*\alpha=d(w^*f)$，故最后一个积分为零。
\end{proof}

\stfhIndex{I078}{wei^xing^kong^jian}{位形空间（configuration space）}
在正合辛流形 $(M,d\alpha)$ 上，可以考虑作用量泛函
\[
 \mathcal A_H(\gamma)=-\int_\gamma\alpha
 -\int_0^1H(t,\gamma(t))\,dt
\]
及其第一变分
\begin{equation}
\begin{split}
 d\mathcal A_H(\gamma)(\xi)
 ={}&\int_0^1\omega\bigl(\xi,\dot\gamma-X_H(t,\gamma)\bigr)\,dt\\
 &+\langle\alpha(\gamma(0)),\xi(\gamma(0))\rangle
  -\langle\alpha(\gamma(1)),\xi(\gamma(1))\rangle.
\end{split}
\label{eq:exact-action-first-variation}
\end{equation}
这无非是把\eqref{eq:conormal-first-variation}中的 $\theta$ 换成 $-\alpha$。周期边界条件仍是使边界贡献消失的自然边界条件。

\begin{question}\label{ques:other-natural-boundary}
这个泛函 $\mathcal A_H$ 是否还有别的自然边界条件？
\end{question}

\begin{remark}\label{rem:d-branes-natural-boundary}
这类问题是物理学家为给具有适当\emph{物理拉格朗日量}的给定系统正确定义\emph{运动方程}而经常提出的问题的原型。在\emph{弦论}中，作为给定系统自然边界条件出现的对象大体对应于\emph{D-膜}的概念。大约自 1996 年物理学家在开弦理论中引入 D-膜概念以来，它一直在当代弦论中起基础作用。
\end{remark}

下面尝试回答上述问题。首先注意，当我们对某些子流形 $L_0,L_1$ 施加边界条件
\[
 \gamma(0)\in L_0,\qquad\gamma(1)\in L_1
\]
时，\eqref{eq:exact-action-first-variation}中的两个边界项分别就是
\[
 i_0^*\alpha(\xi(0)),\qquad i_1^*\alpha(\xi(1)).
\]
现在假设 $L_0,L_1$ 是正合拉格朗日子流形，从而对某些函数
$f_0:L_0\to\mathbb R$ 与 $f_1:L_1\to\mathbb R$ 有
\[
 i_0^*\alpha=df_0,\qquad i_1^*\alpha=df_1.
\]
在此情形下，可以在原作用量泛函 $\mathcal A_H$ 上添加由 $f_0,f_1$ 给出的项，考虑修正后的作用量泛函
$\mathcal A_{H;(f_0,f_1)}:\mathcal P(L_0,L_1)\to\mathbb R$：
\begin{equation}
 \mathcal A_{H;(f_0,f_1)}(\gamma)
 =\mathcal A_H(\gamma)+f_1(\gamma(1))-f_0(\gamma(0)).
 \label{eq:modified-exact-action}
\end{equation}
于是这个修正作用量泛函在 $\mathcal P(L_0,L_1)$ 上的第一变分成为
\[
 d\mathcal A_{H;(f_0,f_1)}(\gamma)(\xi)
 =\int_0^1\omega\bigl(\dot\gamma-X_H(t,\gamma(t))\bigr)\,dt,
\]
从而修正作用量泛函的运动方程为
\[
 \dot z=X_H(z),\qquad z(0)\in L_0,\quad z(1)\in L_1,
\]
这正是我们所要的。由此引出塞德尔 (Se03a) 给出的\emph{正合拉格朗日膜}定义，其中还要给 $L$ 添加一些额外装饰。

\begin{definition}[正合拉格朗日膜]\label{def:exact-lagrangian-brane}
考虑满足 $\omega=d\alpha$ 的 $(M,\omega)$。设
$i:L\hookrightarrow(M,\omega)$ 是正合拉格朗日子流形。若配对 $(L,f)$ 满足
$df=i^*\alpha$，则称它为\emph{正合拉格朗日膜}。
\end{definition}

下面用类似余切丛情形的计算，证明正合拉格朗日膜的一个有趣公式。

\begin{proposition}\label{prop:hamiltonian-flow-exact-brane}
设 $(M,d\alpha)$ 是正合辛流形，$(L,g)$ 是正合拉格朗日膜，包含映射为
$i:L\hookrightarrow M$。哈密顿流 $\phi_H^t$ 诱导一族光滑的正合拉格朗日膜 $(L_t,f_t)$；其中 $L_t=i_t(L)$、$i_t:L\to M$、$f_t:L\to\mathbb R$，并且
\begin{equation}
 i_t=\phi_H^t\circ i,\qquad
 f_t=g+\int_0^t\bigl(H_s+\alpha(X_{H_s})\bigr)\circ i_s\,ds.
 \label{eq:hamiltonian-deformed-brane-phase}
\end{equation}
\end{proposition}

\begin{proof}
由 $f_t$ 的定义，$g=f_0$，所以 $i^*\alpha=df_0$。须证明对所有 $t\in[0,1]$，
\begin{equation}
 i_t^*\alpha=df_t.
 \label{eq:deformed-brane-exactness}
\end{equation}
为此，再次计算
\[
 i_t^*\alpha=(\phi_H^t\circ i)^*\alpha=i^*(\phi_H^t)^*\alpha
\]
的导数。简单计算给出
\begin{equation}
\begin{aligned}
 \frac d{dt}(\phi_H^t)^*\alpha
 &=(\phi_H^t)^*\mathcal L_{X_{H_t}}\alpha\\
 &=(\phi_H^t)^*\bigl(X_{H_t}\mathbin{\lrcorner}d\alpha\bigr)
   +d\bigl((\phi_H^t)^*(X_{H_t}\mathbin{\lrcorner}\alpha)\bigr)\\
 &=d\bigl(H_t\circ\phi_H^t+\alpha(X_{H_t})\circ\phi_H^t\bigr).
\end{aligned}
\label{eq:liouville-form-under-hamiltonian-flow}
\end{equation}
因此，令 $i_0^*\alpha=dg$，便有
\[
\begin{aligned}
 i_t^*\alpha
 &=dg+\int_0^t\frac d{ds}i_s^*\alpha\,ds\\
 &=d\left(g+\int_0^t
   \bigl[H_s\circ\phi_H^s+\alpha(X_{H_s})\circ\phi_H^s\bigr]\circ i\,ds\right)\\
 &=d\left(g+\int_0^t\bigl[H_s+\alpha(X_{H_s})\bigr]\circ i_s\,ds\right).
\end{aligned}
\]
这就由上面对 $f_t$ 的定义证明了\eqref{eq:deformed-brane-exactness}。
\end{proof}

\begin{remark}\label{rem:brane-phase-monodromy}
需要指出，即使 $L_1=L_0$，函数 $g_1$ 也可能不同于 $g_0=g$。正合拉格朗日膜的相位函数所呈现的这种单值性现象，在塞德尔对正合拉格朗日膜的深谷范畴的研究 (Se03a) 中产生了一个有趣结果。关于这一问题的进一步讨论参见 (Se03a)。
\end{remark}
